\begin{abstract}

\todo{DO NOT REMOVE: Uncomment the todo relax macro and make sure this line disappears.}


Recent studies have revealed 
that self-sustaining cascading failures 
in distributed systems 
frequently lead to widespread outages, 
which are challenging to contain and recover from. 
Existing failure detection techniques 
struggle to expose such failures prior to deployment, 
as they typically require a complex combination of specific conditions to be triggered. 
This challenge stems from the inherent nature of cascading failures, 
as they typically involve a sequence of fault propagations, 
each activated by distinct 
conditions. 

This paper presents \tool{}, 
a fault injection 
framework 
to expose self-sustaining
cascading failures in distributed systems. 
\tool{} 
uses the novel idea of \textit{causal stitching}, which causally links multiple single-fault injections 
in different test workloads
to simulate complex fault propagation chains.
To identify propagation chains between faults,
\tool{} designs 
a \textit{counterfactual} causality analysis of fault propagations 
-- \textit{fault causality analysis} (FCA): 
FCA compares the execution trace of a fault injection run 
with its corresponding profile run 
(i.e., running the same test without the injection) 
and identifies any additional faults triggered,
which are considered to have a causal relationship with the
injected fault. 

To address the large search space of fault and workload combinations,  
\tool{} employs a \textit{three-phase allocation (3PA)} protocol of test budget 
that prioritizes faults with unique and diverse causal consequences, 
thereby increasing the likelihood of uncovering conditional fault propagations. 
Furthermore, 
to avoid incorrectly connecting fault propagations from workloads with incompatible conditions, 
\tool{} performs a \textit{local compatibility check} 
that \textit{approximately} 
checks the compatibility of 
the path constraints associated 
with connected fault propagations 
with low overhead. 


\tool{} has detected \NumOfNewVC bugs that 
resulted in self-sustaining 
cascading failures 
in \NumOfSysTested widely deployed distributed systems, 
\NumOfConfirmedVC of which have been confirmed 
with \NumOfFixedVC fixed. 

\end{abstract}